% ===========================================================================
%  05 — 实现
%
%  声明: C-ARCH-01, C-ARCH-02, C-ARCH-03, C-VLM-01, C-VLM-02,
%        C-TELEM-01, C-RAG-01, C-PROC-01, C-RUNTIME-01
% ===========================================================================
\chapter{实现}\label{ch:implementation}

第~\ref{ch:system-design}章和第~\ref{ch:methodology}章描述了基于证据的教学系统的架构与方法。本章阐述这些思想如何被实现为一个可运行的研究原型系统。论述重点是有选择性的：仅当实现细节用于执行证据边界、保持清洁/六边形架构分离、使帮助循环可审计，或支撑第~\ref{ch:evaluation}章所用的评估产物时，才予以包含。

因此，本章的实现问题并非简单地回答存在哪些模块，而是回答：在生成的教学帮助到达学习者之前，在何处进行了权限检查。答案分布在多个层面：从过程包派生的 Schema、证据包、响应映射、Harness 验证、有边界的叠加层执行、降级逻辑，以及追加式事件日志。不改变上述论证的操作性细节（如完整的命令参考和完整部署拓扑）置于附录~\ref{ch:appendix}中。

% ===========================================================================
\section{运行时边界与配置}\label{sec:impl-stack}

代码库围绕清洁/六边形架构边界组织。领域逻辑位于 \texttt{core/}；稳定接口在 \texttt{ports/} 中声明；外部系统通过 \texttt{adapters/} 连接；任务相关的程序知识存储在 F/A-18C 过程包中。这种布局是一种实现决策，而非仅仅是代码风格偏好。它将仿真器传输、模型提供方、视觉采集、检索、叠加层分发和实验导出等外部关注与决定什么证据可采信的领域对象隔离开来。

Python 项目目标为 Python~3.10+，使用 \texttt{pyproject.toml} 中声明的依赖集。运行时依赖少量通用库：\texttt{httpx} 用于 HTTP 模型调用，\texttt{PyYAML} 用于声明式配置，\texttt{jsonschema} 用于契约验证，\texttt{Pillow} 用于 VLM 路径中的图像处理。测试使用 \texttt{pytest} 实现。这些技术栈选择之所以重要，主要是因为它们支持可检查的基于文件的配置和可复现的验证。

模型访问在设计中保持提供方无关。文本帮助路径使用 OpenAI 兼容适配器，可配置用于 vLLM、llama.cpp 或 TGI 等服务器。通过 Ollama 适配器提供本地兼容性，确定性 stub 则支持测试。同样的边界也用于多模态 VLM 路径：运行时将视觉事实提取视为返回结构化观测的适配器，而非程序决策的权威来源。

DCS 端集成扮演着同样狭窄的角色。Lua 脚本和 sidecar 组件导出 DCS-BIOS 状态、触发复合面板截图，并应用叠加层高亮指令。它们不决定当前过程步骤、不生成教学文本，也不绕过 Python 端的验证链。这使得仿真器集成保持为 I/O 适配器，而过程包、证据包和 Harness 仍然是教学决策的权威来源。

% ===========================================================================
\section{证据根基执行}\label{sec:impl-evidence}

证据根基通过多个实现关卡来执行，而非仅仅通过提示词措辞。模型可以提出诊断、下一步建议、解释和叠加层目标，但该提议必须通过从过程包生成的契约，以及针对当前帮助循环中可用证据的检查。

\paragraph{从过程包派生的 Schema。}
函数 \texttt{build\_help\_response\_schema} 在运行时生成 \texttt{HelpResponse} JSON Schema。步骤标识符来自步骤注册表，叠加层目标来自 UI 映射，错误类别来自分类法。Schema 要求每个叠加层目标至少有一项相应的证据项。证据项必须使用下列运行时根基类型之一：\texttt{var}、\texttt{gate}、\texttt{rag}、\texttt{delta} 或 \texttt{visual}。这使得输出契约针对具体过程，而非一个通用的语言模型格式。

\paragraph{证据包构建。}
帮助提示词是由结构化上下文而非原始日志片段构建的。\texttt{build\_evidence\_packet} 将遥测证据、门控证据、最近操作、视觉事实、确定性候选和已知冲突聚合成一个紧凑的证据包。\texttt{build\_help\_prompt\_result} 然后将此证据包、允许的证据引用、检索到的片段以及叠加层目标策略纳入面向模型的提示词中。因此，证据引用可以在之后针对展示给模型的同一上下文进行检查。

\paragraph{模型输出映射与 Harness 验证。}
响应路径故意保持狭窄。JSON 提取会移除代码围栏并读取第一个有效对象，但不会将任意格式错误的文本修复为新的答案。解析后的对象先对照动态 Schema 进行验证，然后传递给响应映射器。映射器检查叠加层证据引用是否存在于帮助循环上下文中，对照白名单验证叠加层目标，在元数据中记录映射失败，并应用配置的叠加层目标数量限制。

Harness 层添加了第二个实现边界。提示词暴露了用于候选决策的 \texttt{HarnessDecision} 契约，而公共运行时输出仍然通过稳定的 \texttt{HelpResponse} 路径映射。\texttt{plan\_harness\_action} 然后检查所选步骤、候选证据、允许的叠加层目标、证据引用和最终状态一致性。它可以接受模型提议、将其修复为更安全的目标、强制仅文本指导，或拒绝叠加层输出。这是架构原则"LLM 不是最终权威"的具体实现。

\paragraph{有边界的叠加层执行。}
最终执行器仅接受类型为 \texttt{overlay} 的操作。它拒绝控制注入、按键仿真、空目标、无法映射的目标、超出 UI 白名单的目标，以及需要证据但缺少证据引用的叠加层。当前运行时支持可配置的多目标指导。超过配置限制的目标将被丢弃并记录，而非静默执行。当 DCS 发送端发出确认消息时，请求、应用和失败的叠加层事件会被规范化到与帮助循环其余部分相同的 JSONL 事件流中。

% ===========================================================================
\section{确定性降级}\label{sec:impl-fallback}

降级被实现为一条保守的退化路径。当模型不可用、JSON 提取失败、Schema 验证失败、响应映射拒绝提议的操作，或 Harness 无法生成有根基的叠加层计划时，系统会使用降级路径。降级逻辑检查当前活动过程步骤、未满足的门控条件、最近操作、UI 目标绑定和证据可用性。

重要的是区分硬阻塞和软阻塞。硬阻塞意味着观测到的证据显示某个必要条件为假。软阻塞意味着所需的证据源缺失或不可观测。对于硬阻塞，降级可以给出纠正性指导；但对于软阻塞，当必要证据缺失时，降级路径避免声称学习者未能完成某个步骤。叠加层输出被抑制，除非本地过程规则识别出唯一有效的目标，或当前配置明确允许有边界的多目标计划。

这条降级路径同样维护了论文声明的边界。帮助循环可以在没有模型或没有视觉证据的情况下继续，但由此产生的指导会更加保守：它依赖观测到的遥测数据、过程门控和确定性步骤提示，而非编造的视觉或教学解读。

% ===========================================================================
\section{VLM 视觉事实提取运行时}\label{sec:impl-vlm}

VLM 运行时被实现为视觉事实提取器，而非教学导师。当前过程包将 S08、S15、S18 和 S19 标记为视觉优先步骤。对于这些步骤，实时运行时可以在组装最终帮助上下文之前请求复合面板截图。对于其他步骤，过期的视觉事实不允许覆盖更新鲜的遥测或过程门控证据。

截图请求被发送到 DCS 端的视觉 sidecar，后者将左 DDI、AMPCD 和右 DDI 渲染为复合图像。Python 端的缓冲视觉会话保留最近的观测记录，并选择帮助触发时刻附近的帧。在实时模式下，默认同步窗口为 250\,ms；回放使用更窄的确定性窗口。所选帧的元数据记录了帧标识符、截图时间、同步状态、与触发时刻的偏差以及过期帧标志。

选定的图像帧被传递给视觉事实提取器。提取器根据当前的 13 项事实本体、复合布局和事实特定的决策边界构建提示词。模型响应被解析为 JSON，经清理后对照一个其事实标识符来自视觉事实包配置的 Schema 进行验证。未知的事实 ID、无效状态、不完整的事实对象和非 Schema 字段将被忽略或记录在元数据中。被接受的事实仅限于 \texttt{seen}、\texttt{not\_seen} 和 \texttt{uncertain} 三种状态。

经过验证的视觉事实按照本体的粘性和过期策略合并到运行时视觉快照中。这一实现细节很重要，因为某些座舱显示状态是瞬态的：某项事实可能需要短暂保持可用以满足过程门控或 Harness 决策，而过期事实必须在误导后续步骤之前失效。如果多模态提取被禁用或 VLM 调用失败，帮助循环仅依靠遥测、检索和降级证据继续。

% ===========================================================================
\section{回放与日志基础设施}\label{sec:impl-replay}

该实现将日志视为证据链的一部分。运行时事件被写入 JSONL 存储，并在每次追加后立即刷新。因此，一个帮助循环会留下可审计的轨迹，包括触发元数据、遥测上下文、视觉帧选择、提示词构建元数据、模型和映射状态、Harness 决策、叠加层执行和失败原因。

回放使用与实时操作相同的边界准则。预录的 DCS-BIOS 捕获可以通过运行时回放，并可选择附带视觉帧目录，而过程包、UI 映射、遥测映射、源策略和 Harness 规格仍然是显式输入。\texttt{fa18c\_startup\_v04} 回放套件组合了五个场景案例和实时 fixture 回归条目，覆盖了 VLM 过期证据、错误目标阻止、移动控件稳定、已完成的完成条件以及交互文本语义等情况。

这些日志和回放套件并不证明学习效果。它们的技术作用是：使帮助循环决策可复现，检查叠加层为何被接受、修复、丢弃或拒绝，并为第~\ref{ch:evaluation}章报告的 RQ1 就绪性证据提供支撑。

% ===========================================================================
\section{实时运行时编排}\label{sec:impl-live}

实时运行时通过一条共享的帮助循环路径连接各实现组件。帮助触发首先从 DCS-BIOS 帧创建遥测快照。遥测管线将原始观测富化为运行时变量和最近增量。如果当前步骤需要视觉确认，则发出截图请求，缓冲视觉会话选择同步的帧。过程引擎和门控评估器生成确定性候选和缺失条件诊断。这些观测被组装为证据包，检索添加源受控的过程知识片段，提示词构建器构造模型请求。

文本模型返回结构化响应而非自由形式的指令。该响应被解析、Schema 验证、映射为运行时教学响应、经 Harness 检查，并转换为最多不超过配置数量的叠加层操作。叠加层执行器仅将经过验证的叠加层意图分发到 DCS，并记录确认或失败。完整的帮助循环（包括降级和修复元数据）被追加到 JSONL 日志中。

在可能的情况下，回放和回归工具也使用相同的编排路径。共享路径对论文而言很重要，因为评估并非基于一个独立的教学系统模拟实现。约束实时教学的同一套契约也产出了用于评估技术就绪性的日志、回放输出和质量门控。
